\section{Introduction}\label{sec:intro}

\PARstart{R}{unning} an old and a new version of a program on the same inputs and comparing what they do is one of the oldest and most widely used ways of checking a change~\cite{evans2007differential}, an application of differential testing~\cite{mckeeman1998differential}. When the two versions disagree, the result is actionable: the differing input is a witness that can be replayed, minimised and inspected. When they agree on every input, the result is much harder to interpret. Evans and Savoia already observed that an empty difference set is ambiguous~\cite{evans2007differential}: it may mean that the change preserved behaviour, or that the inputs never reached the changed code, reached it in only one redundant way, or produced a difference on a channel that was not observed. Despite this, an empty difference set is routinely read as reassurance, both by developers reviewing a change and by automated pipelines that treat ``no difference'' as a pass.

This paper studies the reliability of that reassurance. We call the verdict issued when a differential comparison observes no divergence \NDOfull{} (\NDO{}), and we ask how often such a verdict is wrong. The question is not whether the two versions are equivalent; that is undecidable in general and is the domain of regression verification~\cite{godlin2009regression,lahiri2012symdiff}. The question is narrower and empirical: given the evidence that supported a particular \NDO{} verdict, how likely is it that a more thorough examination of the same change, within a declared observation boundary and adjudication budget, would have found a divergence? We call this probability the \emph{false-assurance risk} of the verdict.

Several bodies of work bear on the question without answering it directly. Coverage criteria quantify how much of a program the evidence exercised, but coverage correlates only weakly with fault-detection effectiveness once suite size is controlled~\cite{inozemtseva2014coverage,gopinath2014code,kochhar2015code}. Mutation analysis asks whether a passing suite would have detected seeded faults, and is therefore a principled way to evaluate passing outcomes~\cite{jia2011analysis,just2014mutants,papadakis2019mutation}; checked coverage and pseudo-tested methods ask whether executed code is actually checked by an oracle~\cite{schuler2011assessing,niedermayr2016will,veraperez2019comprehensive}. These techniques characterise a test suite for a single program version; they do not attach a calibrated error probability to a specific no-divergence verdict on a specific version pair. Selective classification supplies exactly such an estimand---the error rate among accepted predictions at a given coverage~\cite{elyaniv2010foundations,geifman2017selective}---and conformal methods supply guarantees under exchangeability~\cite{angelopoulos2024conformal,bates2021distribution}, but neither defines what the evidence behind a software verdict is, nor how an estimate should respond when that evidence is degraded or tampered with. Recent work measures how accurate agreement-based ``semantics-preserving'' labels are over a whole dataset~\cite{ayub2026semadiff}; it does not estimate a per-verdict risk or abstain when the evidence is weak.

We investigate whether this gap can be closed by treating a no-divergence verdict as a selective prediction whose reliability is estimated from a structured description of its supporting evidence. The investigation is implemented in \sbg{}, the eighth iteration of the Software Behavior Genome research prototype, and has four parts.

\begin{enumerate}
\item \emph{An explicit estimand.} We define the false-assurance risk $\R = P(F{=}1 \mid \NDO{}\text{ issued})$ relative to a declared observation boundary and evaluation budget, state how unresolved cases enter it, and separate it from equivalence (Section~\ref{sec:problem}).
\item \emph{A change-conditioned evidence profile.} Rather than counting inputs, the profile records whether the evidence activated the changed code, how many independent and reproducible observations reached it, which observation channels were present, and whether provenance, integrity and environment checks held (Section~\ref{sec:approach}). A calibrated model over this profile, which we call change-conditioned evidence calibration (\cec{}), maps it to an estimated risk $\Rhat$ and an abstention decision.
\item \emph{An Evidence Response Signature (\ers{}).} We specify eight properties that a reasonable estimator of $\R$ should satisfy---for example, adding irrelevant or duplicate evidence should not lower the estimated risk, and an integrity violation should prevent a reassuring verdict---and nine controlled evidence perturbations (P1--P9) that test them (Section~\ref{sec:ers}). \ers{} is used in two roles that we keep separate throughout: as a \emph{conformance test} of an estimator and as a source of \emph{predictive features}.
\item \emph{Trusted evidence collection and a matched-coverage evaluation protocol.} Coverage and run records are captured by a controller outside the program under test, standard output is teed by an in-process stub, an optional shadow run re-observes return values, and the evidence is sealed, and the architecture is subjected to adversarial probes (Section~\ref{sec:trust}). All risk comparisons are made at matched coverage, because an estimator can lower observed false-assurance risk simply by issuing fewer verdicts (Section~\ref{sec:design}).
\end{enumerate}

We evaluate these components on a development corpus of 312 executed version pairs. The 258 labelled pairs come from generated (synthetic) programs and from QuixBugs~\cite{lin2017quixbugs}; 46 BugsInPy-derived pairs were executed but remain unlabelled, so they enter the perturbation analyses but not the risk estimates. The results are mixed, and we report them as they are. As a conformance test, \ers{} was useful: an all-feature probe estimator failed the irrelevant-evidence invariance property on 109 of 227 applicable records, exposing a defect in the evidence profile that was then fixed and re-verified (227 of 227 pass), and the property matrix exposed that the predecessor's verdicts ignore integrity failures (155 of 312 records when replayed, 84 when recomputed), which motivated the \sbg{} trust gate that now passes by construction. As a predictor, the richer evidence was not shown to help: at the coverage of the predecessor's sufficiency gate (61 of 258 labelled pairs), the full \cec{} model with \ers{} features had an observed false-assurance risk of 0.213 against 0.230 for coverage-only calibration, a 1.6-percentage-point difference whose project-clustered 95\% confidence interval spans $-6.6$ to $+16.3$ points and whose sign changes across cross-validation replicates. A power analysis shows that the development corpus could only have detected differences of roughly 16 points; detecting the pre-registered 5-point improvement would require on the order of 460 to 1{,}160 issued \NDO{} verdicts.

The paper therefore offers a methodological contribution and a transparent pilot evaluation rather than a claim of improved prediction. Specifically, it contributes:
\begin{itemize}
\item[C1)] a formal false-assurance estimand for no-divergence verdicts and its relationship to selective risk;
\item[C2)] a change-conditioned evidence profile that includes integrity, provenance and environment signals;
\item[C3)] the \ers{} property set and the P1--P9 perturbation protocol, with evidence that the protocol detects estimator defects and discriminates between decision rules;
\item[C4)] a trusted evidence-collection architecture with an explicitly bounded threat model and adversarial results that distinguish detection from prevention;
\item[C5)] a matched-coverage evaluation of \cec{} against rule baselines and ablations, with an inconclusive result and a power analysis that sizes the confirmatory study; and
\item[C6)] a reproducibility package and an audit of the project's own artifacts, including the exclusion of an evaluation artifact that did not meet the evidential standard (Section~\ref{sec:repro}).
\end{itemize}

Equally important is what the paper does not claim. It does not claim that \NDO{} establishes equivalence, that \sbg{} detects every behavioural divergence, that post-hoc evidence perturbations estimate causal effects, that the \ers{} features improve risk prediction, that the evidence collector is secure against arbitrary adversarial programs, or that the results generalise beyond the development corpus. Independent human adjudication of an external corpus has been designed and implemented but has not been carried out.

Section~\ref{sec:background} reviews related work and positions \sbg{} against it. Section~\ref{sec:problem} formalises the problem. Sections~\ref{sec:approach}--\ref{sec:trust} describe the evidence model, \ers{} and the trusted execution architecture. Section~\ref{sec:design} presents the research questions and experimental design, Section~\ref{sec:results} the results, and Sections~\ref{sec:discussion}--\ref{sec:repro} discuss implications, threats to validity and reproducibility.
